% HIKET -- DISCUSSION MEMO: a running collection of findings that are useful for the
% Discussion but not yet strong enough for the storyline.
%
% THIS IS A LIVING DOCUMENT. Items are added, promoted, demoted and DELETED. It is not
% a record of everything we have done -- doublechecks/ and the git history are that.
% It holds only what might earn a sentence in the Discussion, together with the honest
% reason it is not there yet.
%
% HOW TO MAINTAIN IT (see also the "Rules" section in the body):
%   - one \item per finding, using the \finding environment, newest at the top of its
%     status group;
%   - when an item becomes strong enough, MOVE it to the storyline/manuscript and
%     delete it here, leaving one line in the "Graduated" table;
%   - when an item is falsified, delete it and leave one line in the "Retired" table
%     with what killed it. Do not keep dead items in the body.
\documentclass[11pt]{article}
\usepackage[utf8]{inputenc}\usepackage[T1]{fontenc}
\usepackage[margin=2.4cm]{geometry}
\usepackage{graphicx}\usepackage{amsmath}\usepackage{amssymb}\usepackage{xcolor}
\usepackage{booktabs}\usepackage{longtable}
\usepackage{float}\usepackage[hidelinks]{hyperref}
\usepackage{caption}\captionsetup{font=small,labelfont=bf}
\usepackage[most]{tcolorbox}
\usepackage{enumitem}

% --- status marks -------------------------------------------------------------------
\definecolor{stgood}{HTML}{2E7D4F}\definecolor{stmid}{HTML}{B8860B}\definecolor{stweak}{HTML}{8B5E3C}
\newcommand{\stA}{\textcolor{stgood}{$\bullet\bullet\bullet$}}  % measured, robust, one step from the storyline
\newcommand{\stB}{\textcolor{stmid}{$\bullet\bullet\circ$}}     % measured but confounded or single-source
\newcommand{\stC}{\textcolor{stweak}{$\bullet\circ\circ$}}      % suggestive; would need new data or a new test

% --- the finding environment --------------------------------------------------------
\definecolor{fbg}{HTML}{F4F6F8}\definecolor{frule}{HTML}{5B6B7A}
\newtcolorbox{finding}[1]{colback=fbg, colframe=frule, boxrule=0pt, leftrule=3pt,
  arc=1pt, left=8pt, right=8pt, top=5pt, bottom=5pt, fontupper=\small,
  title={#1}, coltitle=frule, fonttitle=\small\bfseries,
  attach title to upper=\par, before skip=10pt, after skip=10pt}

\definecolor{cavbg}{HTML}{FDF3E7}\definecolor{cavrule}{HTML}{C87A2F}
\newtcolorbox{caveat}{colback=cavbg, colframe=cavrule, boxrule=0pt, leftrule=3pt,
  arc=1pt, left=8pt, right=8pt, top=5pt, bottom=5pt, fontupper=\small,
  before skip=8pt, after skip=8pt}

\newcommand{\sinit}{\sigma_{\mathrm{init}}}
\newcommand{\sinput}{\sigma_{\mathrm{input}}}

\title{HIKET --- discussion memo\\
  \large Findings that may serve the Discussion, and are not yet strong enough for the storyline}
\author{}
\date{Started 12 August 2026 --- living document, revised continuously}

\begin{document}
\maketitle

% =====================================================================================
\section*{What this is}
\addcontentsline{toc}{section}{What this is}

A holding place. Analyses accumulate faster than the storyline can absorb them, and the
ones that do not fit are either lost or, worse, quietly promoted without the caveat that
kept them out. This document keeps them with their caveat attached.

\medskip
\noindent\textbf{Admission rule.} An item belongs here if it could plausibly earn a
sentence in the Discussion \emph{and} there is a stated reason it cannot yet. If there is
no such reason, it belongs in the manuscript. If there is no plausible sentence, it
belongs in \texttt{doublechecks/} and nowhere else.

\medskip
\noindent\textbf{Status marks.}\quad
\stA~measured and robust, one step from the storyline \quad
\stB~measured but confounded, or resting on a single source \quad
\stC~suggestive; needs new data or a new test

\medskip
\noindent\textbf{Maintenance.} Items are \emph{deleted} when they graduate or die, leaving
one line in the tables at the end. This document should not grow monotonically; if it
does, it has become a diary rather than a memo.

% =====================================================================================
\section{Open items}

% -------------------------------------------------------------------------------------
\begin{finding}{\stA\ \ D4 --- The litter product's between-plot gradient is eighteen-fold; the soil's is $1.2$}
\textbf{Claim.} Across the observed basal-area range the litter input product varies by a factor of
$\sim$18 between plots. Observed soil carbon varies by $1.2$. The models transmit the input gradient
faithfully, so the plot-scale mismatch is a property of the \emph{input product} and of the
\emph{soil}, not of the decomposition models.
\end{finding}

\noindent\textbf{Measured} (arm B; \texttt{doublechecks/basal\_area\_sign.R},
\texttt{woody\_compartment\_test.R}), as $d\log(\cdot)/d\mathrm{BA}$ and as span over
$40$~m$^2$\,ha$^{-1}$:

\begin{center}\small
\begin{tabular}{lrr}
\toprule
 & slope & span\\
\midrule
litter input $J_{t_0}$, matched in time to BA$_{85}$ & $+0.072$ & $\times 18.1$\\
modelled stock (transient-init endpoint, six models) & $+0.073$ to $+0.076$ & $\times 18.6$--$20.5$\\
model predictions after the forward run & $+0.044$ to $+0.049$ & $\times 6.3$\\
observations & $+0.0044$ & $\times 1.19$\\
observations, adjusted for region and soil type & $+0.0018$ ($p=0.12$) & $\times 1.08$\\
\bottomrule
\end{tabular}
\end{center}

\noindent The residual falls with basal area in all six models ($p$ from $10^{-16}$ to $10^{-111}$,
holdout, calibration and pooled alike). The models neither invent nor amplify the gradient --- the
forward run in fact \emph{damps} it.

\begin{caveat}
\textbf{Two of my own errors are corrected here; do not reintroduce them.}
\textbf{(1)} An earlier version of this item put $d\log J/d\mathrm{BA}$ at $0.016$ and concluded the
models amplify the input gradient $\sim$3$\times$. That number came from averaging the annual
\texttt{inputs\_by\_plot} series over all years while regressing on \emph{1985} basal area. Matched in
time --- \texttt{J\_t0\_mean}, the flux the initialiser actually uses --- the input slope is $0.072$
and there is \textbf{no amplification}.
\textbf{(2)} The consequent hypothesis, that woody-derived carbon inflates high-biomass plots because
Yasso's stock is DOM${+}$SOM while our target is a soil measurement, was \textbf{tested and refuted}:
zeroing woody inputs moves the slope $6\%$ (Yasso15 $0.0755 \to 0.0710$) against a $15\times$ gap.
The free control is decisive --- SP1/TP2/TP3 carry a single scalar flux with no composition
information at all and give the \emph{same} slope, so a mechanism absent from three of the six cannot
explain a gradient present in all six.
\end{caveat}

\medskip\noindent\textbf{What survives, and it is more useful than what died.} An eighteen-fold
between-plot input gradient is not credible on its face. A stand-keyed litter product gives a freshly
cut plot almost no input, and a freshly cut plot does not stop receiving carbon --- ground vegetation,
dying roots and logging residues continue. So the plot-scale failure is evidence for the \emph{same}
missing flux as the input argument, reached from the opposite end, and it points specifically at
\textbf{understorey and early-succession vegetation}, the two components that argument is weakest on.
\textbf{It is also independent evidence:} the national-scale argument comes from a level over time,
this one from a gradient across space.

\begin{caveat}
\textbf{The alternative remains live and is not exclusive.} Plot-level soil carbon may simply not be
set by current litter input, because each plot sits at a different point on its own disturbance
trajectory and a single global initial state cannot express that --- this paper's thesis one level
down. The observed gradient being far \emph{smaller} than the input gradient is consistent with both.
\ \textbf{TEST:} the discriminating test would be residuals against time since disturbance, which our plots
support.
\end{caveat}

\begin{caveat}
\textbf{The compartment question itself is not dead, only demoted.} Whether our soil-only target
should be compared with a DOM${+}$SOM model total is still a legitimate question about the
\emph{level} --- it is simply not the explanation for the plot-scale gradient. Zeroing woody inputs
removes $\sim$22\% of the input in the Yasso models, which is a level effect worth knowing when the
$\sinput$ argument is written, since a compartment mismatch and a missing flux push the apparent input
requirement in the same direction. \textbf{Not the same question as depth}, which was checked and is
fine (memory \texttt{yasso-depth-basis-matches}).
\end{caveat}

% -------------------------------------------------------------------------------------
\begin{finding}{\stB\ \ D1 --- The 2006--2024 model--observation disagreement is confined to the mineral soil}
\textbf{Claim.} The six models do not fail to reproduce the recent SOC trend in general.
They reproduce the \emph{organic layer}, and disagree with the \emph{mineral soil}.
\end{finding}

\noindent Over 2006--2024, on the 429 plots measured in both campaigns. The model trajectories are
\emph{unweighted} plot means, so the unweighted column is the one that may be compared with them;
the region-weighted column is the nationally representative figure and is shown so the two bases
are never mixed.

\begin{center}\small
\begin{tabular}{lrr}
\toprule
compartment & unweighted & region-weighted\\
 & \multicolumn{2}{c}{Mg\,C\,ha$^{-1}$\,yr$^{-1}$}\\
\midrule
organic layer & $-0.094$ & $-0.082$\\
mineral 0--40\,cm & $+0.232$ & $+0.181$\\
whole measured profile & $+0.138$ & $+0.098$\\
\midrule
the six models (unweighted) & \multicolumn{2}{c}{$-0.02$ to $-0.43$}\\
\bottomrule
\end{tabular}
\end{center}

\begin{caveat}
\textbf{Basis warning, and it bites elsewhere.} The observed trend quoted as $+0.209$ in
\texttt{HIKET\_next\_session} and \texttt{M\&M\_parameterization\_working\_document} is the
\emph{unweighted} paired 2006--2024 rate. The region-weighted equivalent is $+0.139$, and LUKE's
official weighted figure is $+0.111$. All three are defensible; mixing them is not. Model
comparisons must use unweighted; anything described as national must use weighted.
\end{caveat}

\noindent The organic layer ($-0.094$ unweighted) falls \emph{inside} the model range. The entire discrepancy sits
in the mineral soil --- which is independently the least reliable compartment: samples are
cut from the wall of a spade-dug pit rather than taken volumetrically
(Kramarenko 2012, \S4.4), proportional bias appears between \emph{every} campaign pair in
both mineral layers while the humus layer shows none, and $\sim$75\% of the first
campaign's mineral carbon concentrations were predicted from loss-on-ignition rather than
measured.

\medskip
\noindent Supporting arithmetic: the organic decline is also of the magnitude a falling
litter input predicts. For a first-order pool a fractional input decline $\delta$
propagates as $\delta\,(1-e^{-kt})$; with an organic-layer turnover of 30--40\,yr over an
18-year window that factor is $\approx 0.4$, so the $\sim$13\% input decline in the Tupek
series predicts $\sim$5\% stock loss against $-6.6$\% observed. A 5-year-turnover pool
would have predicted $\sim$13\%, which is not what happened.

\begin{caveat}
\textbf{Why this is not in the storyline.}
(i) The organic layer is still a \emph{layer}, defined by a field boundary, and that
boundary moved: $\Delta$organic against $\Delta$mineral\,0--20 gives slope $-0.25$
($r=-0.16$, $p=0.0007$). The movement accounts for only $\sim$25\% of the organic loss, so
$\sim$75\% survives --- but the compartment is not boundary-proof. Only the whole-profile
total is, and that shows $+0.20\pm0.11$\,kg\,m$^{-2}$, \emph{not significant}.
(ii) The organic total is the sum of a litter layer that rose 27\% and a humus layer that
fell 12\%, and neither component behaves sensibly on its own (item D2). The aggregate may
be right for the wrong reasons.
(iii) The model range quoted is from the $\sigma=0.72$ run and predates both the litter and
the dating corrections to the target.
\end{caveat}

\noindent\textbf{What would promote it.} Re-running the comparison after the 1985 litter and
sampling-year corrections land; and a model-side check that the models' \emph{own} organic
pool, not just their total, matches. \textbf{What would kill it.} Evidence that the
organic/mineral boundary moved much more than the $-0.25$ slope implies.
\emph{Code:} \texttt{doublechecks/organic\_mineral\_boundary.R}. \emph{Added 2026-08-12.}

% -------------------------------------------------------------------------------------
\begin{finding}{\stC\ \ D2 --- The measured litter layer does not behave like a litter pool}
\textbf{Claim.} The litter layer (LM) grew 27\% between 2006 and 2024 while its input fell,
and no plot-level mechanism explains it.
\end{finding}

\noindent Aggregate: input $-4.9$\% (2004--08 $\to$ 2020--24), litter layer $+27.7$\%, humus
$-11.8$\%. Per plot, neither candidate mechanism holds: growth tracking input gives
$r=+0.09$ ($p=0.07$), and reclassification out of the humus gives $r=-0.04$ ($p=0.45$,
slope $-0.01$). Cross-sectionally the litter layer has no relation to modelled litter input
in either campaign ($r=+0.03$, $+0.01$).

\medskip
\noindent The surviving explanation is a \emph{uniform campaign-level} change in what was
called litter in 2024: plot-level correlations are blind to an offset applied equally to
every plot, which is exactly the signature observed.

\begin{caveat}
The null cross-sectional correlation cannot separate ``the litter layer is not behaving
like a litter pool'' from ``the Tupek series carries little plot-level signal''. The latter
is independently known (litter does not order SOC across plots; $R^2\approx0$). So the
real-growth hypothesis is closer to \emph{untestable with this predictor} than rejected.
This must not be reported as evidence that the litter layer is spurious.
\end{caveat}

\noindent\textbf{What would promote it.} An independent plot-level litter predictor with
demonstrated skill; or confirmation from LUKE that the 2024 litter sampling instruction
differed. \emph{Code:} \texttt{doublechecks/litter\_layer\_vs\_input.R}. \emph{Added 2026-08-12.}

% -------------------------------------------------------------------------------------
\begin{finding}{\stA\ \ D3 --- Plot-level SOC change is not resolvable between campaigns}
\textbf{Claim.} The 95\% limits of agreement between any two campaigns are of the same size
as the stock being measured.
\end{finding}

\noindent Bland--Altman, whole profile: $-30.5$ to $+43.6$\,Mg\,C\,ha$^{-1}$ for 1985--2006 and
$-39.0$ to $+42.9$ for 2006--2024 --- about 115\% of the mean stock in both cases. For the
mineral 20--40\,cm layer it is $\sim$220\%. This is a model-free statement: it uses only two
measurements of the same plots.

\medskip
\noindent It is the quantitative form of ``the posteriors are overconfident'', and it needs
no model to establish, which makes it the cleanest available support for treating
plot-years as far less informative than their count suggests.

\begin{caveat}
Placement, not strength, is what keeps this out of the storyline: it is an argument about
the error model and the likelihood, and the storyline does not yet have a home for it. It
may belong in the M\&M working document instead of the Discussion.
\end{caveat}

\noindent\emph{Code:} \texttt{Reporting/August2026\_internal\_report/build\_bland\_altman\_campaigns.R};
figures there and manuscript Fig.~S10(c). \emph{Added 2026-08-12.}

% =====================================================================================
\section{Graduated}

\noindent Items promoted into the manuscript. One line each; the content lives there now.

\begin{center}\small
\begin{tabular}{llp{7.6cm}}
\toprule
date & item & where it went\\
\midrule
--- & --- & \emph{(none yet)}\\
\bottomrule
\end{tabular}
\end{center}

% =====================================================================================
\section{Retired}

\noindent Items removed as falsified or superseded. Kept as one line so they are not
rediscovered and re-argued.

\begin{center}\small
\begin{tabular}{llp{7.6cm}}
\toprule
date & item & what killed it\\
\midrule
2026-08-12 & 1985 depth distribution is distorted & The per-plot decay rate puts 1985
($\lambda=0.038$) \emph{between} 2006 (0.031) and 2024 (0.041). Profile shape does not
single out 1985; the campaign difference is a level, not a shape.\\
\addlinespace
2026-08-12 & The subsoil offset is a stoniness-correction artefact & No relation to coarse
fraction ($\rho=-0.08$, $p=0.15$).\\
\addlinespace
2026-08-12 & 1985 is noisier than the later campaigns & Subsoil repeatability $r=0.63$
(1985--2006) vs $0.65$ (2006--2024); proportional bias present in every pair. Whatever is
wrong with 1985 is a level, not scatter.\\
\bottomrule
\end{tabular}
\end{center}

\end{document}
